% Generated draft: requires mathematical and source review.
\documentclass[11pt]{article}
\usepackage[margin=25mm]{geometry}
\usepackage{fontspec}
\setmainfont{Latin Modern Roman}
\usepackage{amsmath,amssymb,mathtools,unicode-math}
\setmathfont{Latin Modern Math}
\usepackage{enumitem,needspace,xurl,hyperref}
\setlength{\parindent}{0pt}
\setlength{\parskip}{5pt}
\setlength{\emergencystretch}{3em}
\newcommand{\N}{\mathbb{N}}
\newcommand{\Z}{\mathbb{Z}}
\newcommand{\Q}{\mathbb{Q}}
\newcommand{\R}{\mathbb{R}}
\newcommand{\C}{\mathbb{C}}
\newcommand{\F}{\mathbb{F}}
\newcommand{\A}{\mathbb{A}}
\newcommand{\PP}{\mathbb{P}}
\newcommand{\E}{\mathbb{E}}
\newcommand{\eps}{\varepsilon}
\newcommand{\dd}{\,\mathrm{d}}
\begin{document}
\section*{UnsolvedMath ID: 30001755}
\textbf{Source problem code:} catalogue entry
\subsection{Definitions and mathematical statement}
Let \(k\in\Z_{\geq 0}\) and let \(T^k=(S^1)^k\), with \(T^0\) the trivial group. A path-connected space \(X\) is called nilpotent if \(\pi_1(X)\) is a nilpotent group and, for every \(n\geq 2\), its natural action on \(\pi_n(X)\) admits a finite filtration by invariant subgroups on whose successive quotients the action is trivial. A continuous action \(T^k\times X\to X\) is almost free if the stabilizer \(\{t\in T^k:t\cdot x=x\}\) is finite for every \(x\in X\). For a finite CW complex, define its total rational Betti number by
\[
b(X;\Q)=\dim_{\Q}H^*(X;\Q)=\sum_{i\geq 0}\dim_{\Q}H^i(X;\Q).
\]
\textbf{Rational toral-rank conjecture (attributed to Halperin in the supplied record).} For every nonnegative integer \(k\) and every nilpotent finite CW complex \(X\), does the existence of an almost free continuous \(T^k\)-action on \(X\) imply
\[
b(X;\Q)\geq 2^k?
\]
Here the definition of nilpotent space uses the standard path-connected convention; the source record does not state connectedness separately. The supplied references have not been independently checked for this draft, and the record does not explain the role of its auxiliary sectional-curvature reference.
\subsection{Short English statement}
Halperin's rational toral-rank conjecture asks whether every almost free action of a \(k\)-dimensional torus on a nilpotent finite CW complex forces its total rational Betti number to be at least \(2^k\). The supplied record labels this problem open; this draft does not independently verify that status.
\subsection{Sources}
\begin{itemize}
\item Local UnsolvedMath registry, supplied complete record for permanent ID \(30001755\), titled \emph{The Rational Toral-Rank Conjecture}, source problem code \texttt{catalogue entry}, in the Oberwolfach Reports Open Problems collection. The record gives proposed year 2011 and identifies no individual proposer; its background attributes the conjecture to Halperin.
\item \emph{Arbeitsgemeinschaft: Rational Homotopy Theory in Mathematics and Physics} (2011), Oberwolfach Reports, pp.\ 993--1045, DOI \texttt{10.4171/owr/2011/18}. Supplied source and literature reference: \url{https://doi.org/10.4171/owr/2011/18}. The collection describes these reports as documenting mathematical workshops and their open-problem discussions.
\item Supplied auxiliary literature reference: Bettiol--Mendes--Töben, ``bounds on sectional curvature,'' \url{https://arxiv.org/abs/1407.6201}. This label is reproduced from the record; no more precise bibliographic description or connection to the conjecture is supplied.
\item Supplied literature assessment dated \texttt{2026-08-21}: status ``open.'' The curation states that the report explicitly records Halperin's toral-rank conjecture as still open and that no general proof or counterexample has been found. This is supplied dated curation, described as literature triage derived from the OWR curation, not an independently verified assessment. Time-sensitive status and cited claims require independent verification.
\item The registry assigns statement status \texttt{corrected\_verified} and states that the mathematical content was checked against the official report, with notation, context, and question form normalized without changing the source-level claim. This is a supplied verification claim, not verification performed for this draft.
\end{itemize}
\end{document}
